\chapter*{Вступ}
\addcontentsline{toc}{chapter}{Вступ}
\markboth{Вступ}{Вступ}

Злоякісний програмний код --- шкідливе програмне забезпечення, віруси, хробаки, троянські програми, вимагачі, руткіти, шпигунський код --- залишається одним із найпоширеніших і найнебезпечніших інструментів кіберзагроз. За останні десятиліття він пройшов шлях від поодиноких експериментальних вірусів до індустрії, що породжує сотні тисяч нових зразків щодоби, та до високотехнологічної зброї, здатної руйнувати фізичну інфраструктуру й завдавати збитків державного масштабу. Протидія цьому явищу перетворилася з вузькотехнічного завдання антивірусних лабораторій на комплексну наукову та інженерну дисципліну, що поєднує теорію обчислень, математичну статистику, аналіз програм, машинне навчання, мережеві технології, право та організаційне управління.

Україна посідає особливе місце в історії протидії злоякісному коду. Саме тут було вперше застосовано низку найруйнівніших кампаній сучасності --- від атак на енергосистему за допомогою спеціалізованого коду до глобальних деструктивних епідемій, що поширилися через довірені канали постачання програмного забезпечення. Цей гіркий, але безпрецедентно багатий досвід зробив українських фахівців і дослідників носіями знань, цінних для всього світу, і водночас поставив підготовку висококваліфікованих спеціалістів із протидії шкідливому коду в розряд завдань національної безпеки.

Цей навчальний посібник призначено для здобувачів третього (освітньо-наукового) рівня вищої освіти --- майбутніх докторів філософії у сфері кібербезпеки, --- а також для науковців і фахівців-практиків. Його мета --- сформувати цілісне, системне розуміння технологій протидії злоякісному коду: не лише окремих методів і засобів, а й принципів, що їх об'єднують, їхніх теоретичних меж і місця в загальній стратегії захисту. Виклад свідомо зосереджено на \emph{захисних} аспектах --- виявленні, аналізі, протидії та реагуванні; посібник не містить працездатних зразків шкідливого коду й не навчає його створювати, натомість озброює читача знаннями, потрібними для побудови ефективного захисту.

Структурно посібник побудовано за принципом руху від основ до стратегії. Перші розділи закладають фундамент: поняттєвий апарат і таксономію злоякісного коду (розділ~\ref{chap:ch01}), нормативно-правове й стандартизаційне забезпечення (розділ~\ref{chap:ch02}), теоретичні та математичні основи, зокрема принципову нерозв'язність точного виявлення (розділ~\ref{chap:ch03}), моделювання загроз і ланцюги ураження (розділ~\ref{chap:ch04}), внутрішню будову та життєвий цикл шкідливого коду (розділ~\ref{chap:ch05}). Наступні розділи розкривають методи аналізу та виявлення: статичний аналіз і сигнатурне виявлення (розділ~\ref{chap:ch06}), динамічний і поведінковий аналіз (розділ~\ref{chap:ch07}), техніки ухилення й обфускації, яким протистоїть захист (розділ~\ref{chap:ch08}), машинне навчання та статистичне виявлення (розділ~\ref{chap:ch09}).

Далі посібник переходить від окремих методів до їх практичної реалізації та оркестрування: захист кінцевих точок, моніторинг і телеметрія (розділ~\ref{chap:ch10}), інженерія виявлення (розділ~\ref{chap:ch11}), розвідка загроз і атрибуція (розділ~\ref{chap:ch12}), протидія цілеспрямованим сталим загрозам і полювання за загрозами (розділ~\ref{chap:ch13}). Окремі розділи присвячено специфічним середовищам і векторам: протидії шкідливому коду в промислових системах, Інтернеті речей, мобільних і хмарних середовищах (розділ~\ref{chap:ch14}) та атакам на ланцюг постачання програмного забезпечення (розділ~\ref{chap:ch15}). Завершують посібник розділи про реагування на інциденти й цифрову форензику (розділ~\ref{chap:ch16}) та про організаційне управління й стратегію кіберстійкості (розділ~\ref{chap:ch17}), що синтезують увесь матеріал у цілісну систему.

Крізь увесь посібник проходить кілька наскрізних ідей. Перша --- принципова неможливість абсолютного виявлення й абсолютного захисту, що випливає з теорії обчислень і теорії ймовірностей, а отже й необхідність мислити категоріями ризику, а не гарантій. Друга --- перевага стійких, семантично значущих ознак над поверхневими, бо саме вони опираються спробам ухилення. Третя --- ешелонування й каскад методів із зростанням вартості й покриття, у якому дешеві й точні рубежі звільняють дорогі для найскладніших випадків. Четверта --- припущення компрометації й зміщення мети від недопущення атаки до кіберстійкості. П'ята --- нерозривний зв'язок технологій, процесів і людей: найкращі засоби безсилі поза зрілою системою управління.

Кожен розділ супроводжується теоретичним викладом, ілюстраціями, таблицями, формалізованими моделями, практичним кейсом і переліком запитань для самоперевірки, зокрема дослідницьких завдань для самостійної наукової роботи. Наприкінці наведено єдиний список використаних джерел, що охоплює основоположні монографії, міжнародні стандарти, провідні наукові публікації та вітчизняні дослідження. Посібник не замінює, а доповнює практичну роботу з реальними засобами аналізу та захисту й передбачає, що читач поглиблюватиме знання самостійним дослідженням.

Автор висловлює сподівання, що цей посібник стане надійною основою для підготовки дослідників і фахівців, здатних протистояти постійно еволюційним загрозам злоякісного програмного коду, і зробить свій внесок у зміцнення кіберстійкості України.
